% GENERATED FILE. Edit canonical lessons, then run npm run generate:books.
% canonical-source-hash: fnv1a64:46bd3d6ade34dc55
% canonical-lessons: RU-C397-oplachivat, RU-C397-platyozh, RU-C397-besplatnyy, RU-C397-platnyy, RU-C397-vygodnyy

\chapter{Paying for things}
\label{ch:ru-paying-for-things}
\hlchaptermodality{397}

\begin{chapteropening}
\textbf{By the end of this chapter:} \emph{I can say who pays for something, talk about a payment, and say whether something is free, paid for or a good deal.}

Complete the last lesson of chapter 397: I can say who pays for something, talk about a payment, and say whether something is free, paid for or a good deal.
\end{chapteropening}

\section[\texorpdfstring{\ru{оплачивать} (opláchivat)}{оплачивать (opláchivat)}]{\ru{оплачивать} (opláchivat) — to pay for}

\label{lesson:RU-C397-oplachivat}



\begin{quote}
Before the new one: say the Russian for an exchange rate, then the Russian for an exchange.
\end{quote}

\subsection*{You'll want to know: \ru{оплачивать}}
\textbf{\ru{оплачивать}} (\emph{opláchivat}) — \textquotedblleft{}to pay for\textquotedblright{}.

\begin{quote}
\ru{Фирма} \ru{оплачивает} \ru{мне} \ru{обед}.
\end{quote}

\begin{quote}
\emph{Fírma opláchivayet mne obéd.}
\end{quote}

\begin{quote}
The firm pays for my lunch.
\end{quote}

Imperfective; perfective \emph{\ru{оплатить}} (\emph{\ru{оплачу}, \ru{оплатишь}}). It takes a direct object with no \emph{\ru{за}}: \emph{\ru{оплачивать} \ru{счёт}}.

\begin{cousinweb}
\textbf{\ru{оплачивать}} is built from \emph{\ru{о}-} and \emph{\ru{платить}} \textquotedblleft{}to pay\textquotedblright{}, on \emph{\ru{плата}} \textquotedblleft{}a payment\textquotedblright{}, the root that also gives \emph{\ru{зарплата}}.
\end{cousinweb}

\begin{grammarlens}[title={What you've built}]
One more word for telling, explaining and describing at B1.
\end{grammarlens}

\subsection*{Your turn}
\begin{itemize}
\raggedright
  \item \emph{Say it:} \emph{opláchivat}
  \item \emph{Say it:} \emph{opláchivat}, once more
  \item \emph{Say it:} a sentence of your own with \emph{opláchivat}
  \item \emph{Recall:} say the Russian for an exchange rate, then the Russian for an exchange, then say \emph{opláchivat} again
\end{itemize}

\begin{tcolorbox}[breakable,colback=teal!4,colframe=teal!35!black,title={Before you move on}]
What does \ru{оплачивать} mean? (\textquotedblleft{}To pay for\textquotedblright{}.) Say it once more.
\end{tcolorbox}
\section[\texorpdfstring{\ru{платёж} (platyózh)}{платёж (platyózh)}]{\ru{платёж} (platyózh) — a payment}

\label{lesson:RU-C397-platyozh}



\begin{quote}
Before the new one: say the Russian for an exchange, then the Russian for to pay for.
\end{quote}

\subsection*{You'll want to know: \ru{платёж}}
\textbf{\ru{платёж}} (\emph{platyózh}) — \textquotedblleft{}a payment\textquotedblright{}.

\begin{quote}
\ru{Каждый} \ru{месяц} \ru{у} \ru{нас} \ru{большой} \ru{платёж} \ru{за} \ru{квартиру}.
\end{quote}

\begin{quote}
\emph{Kázhdyy mésyats u nas bolshóy platyózh za kvartíru.}
\end{quote}

\begin{quote}
Every month we have a big payment for the flat.
\end{quote}

Masculine; genitive \emph{\ru{платежа}}. \emph{\ru{Ежемесячный} \ru{платёж}} — a monthly payment.

\begin{cousinweb}
\textbf{\ru{платёж}} is built on \emph{\ru{платить}} \textquotedblleft{}to pay\textquotedblright{}, from \emph{\ru{плата}} \textquotedblleft{}a payment\textquotedblright{}, with the suffix \emph{-\ru{ёж}}.
\end{cousinweb}

\begin{grammarlens}[title={What you've built}]
One more word for telling, explaining and describing at B1.
\end{grammarlens}

\subsection*{Your turn}
\begin{itemize}
\raggedright
  \item \emph{Say it:} \emph{platyózh}
  \item \emph{Say it:} \emph{platyózh}, once more
  \item \emph{Say it:} a sentence of your own with \emph{platyózh}
  \item \emph{Recall:} say the Russian for an exchange, then the Russian for to pay for, then say \emph{platyózh} again
\end{itemize}

\begin{tcolorbox}[breakable,colback=teal!4,colframe=teal!35!black,title={Before you move on}]
What does \ru{платёж} mean? (\textquotedblleft{}A payment\textquotedblright{}.) Say it once more.
\end{tcolorbox}
\section[\texorpdfstring{\ru{бесплатный} (besplátnyy)}{бесплатный (besplátnyy)}]{\ru{бесплатный} (besplátnyy) — free (of charge)}

\label{lesson:RU-C397-besplatnyy}



\begin{quote}
Before the new one: say the Russian for to pay for, then the Russian for a payment.
\end{quote}

\subsection*{You'll want to know: \ru{бесплатный}}
\textbf{\ru{бесплатный}} (\emph{besplátnyy}) — \textquotedblleft{}free (of charge)\textquotedblright{}.

\begin{quote}
\ru{Вход} \ru{в} \ru{музей} \ru{бесплатный}.
\end{quote}

\begin{quote}
\emph{Vkhod v muzéy besplátnyy.}
\end{quote}

\begin{quote}
Entry to the museum is free.
\end{quote}

\emph{\ru{Бесплатно}} — for free. Free in the sense of not busy is \emph{\ru{свободный}}.

\begin{cousinweb}
\textbf{\ru{бесплатный}} is built from \emph{\ru{без}} \textquotedblleft{}without\textquotedblright{} and \emph{\ru{плата}} \textquotedblleft{}a payment\textquotedblright{}; the \emph{\ru{з}} turns to \emph{\ru{с}} before \emph{\ru{п}}.
\end{cousinweb}

\begin{grammarlens}[title={What you've built}]
One more word for telling, explaining and describing at B1.
\end{grammarlens}

\subsection*{Your turn}
\begin{itemize}
\raggedright
  \item \emph{Say it:} \emph{besplátnyy}
  \item \emph{Say it:} \emph{besplátnyy}, once more
  \item \emph{Say it:} a sentence of your own with \emph{besplátnyy}
  \item \emph{Recall:} say the Russian for to pay for, then the Russian for a payment, then say \emph{besplátnyy} again
\end{itemize}

\begin{tcolorbox}[breakable,colback=teal!4,colframe=teal!35!black,title={Before you move on}]
What does \ru{бесплатный} mean? (\textquotedblleft{}Free (of charge)\textquotedblright{}.) Say it once more.
\end{tcolorbox}
\section[\texorpdfstring{\ru{платный} (plátnyy)}{платный (plátnyy)}]{\ru{платный} (plátnyy) — paid, that one pays for}

\label{lesson:RU-C397-platnyy}



\begin{quote}
Before the new one: say the Russian for a payment, then the Russian for free.
\end{quote}

\subsection*{You'll want to know: \ru{платный}}
\textbf{\ru{платный}} (\emph{plátnyy}) — \textquotedblleft{}paid, that one pays for\textquotedblright{}.

\begin{quote}
\ru{Здесь} \ru{платная} \ru{стоянка}.
\end{quote}

\begin{quote}
\emph{Zdes plátnaya stoyánka.}
\end{quote}

\begin{quote}
This is a paid car park.
\end{quote}

\emph{\ru{Платная} \ru{дорога}} — a toll road.

\begin{cousinweb}
\textbf{\ru{платный}} is the adjective made from \emph{\ru{плата}} \textquotedblleft{}a payment\textquotedblright{}, the root of \emph{\ru{платить}}.
\end{cousinweb}

\begin{grammarlens}[title={What you've built}]
One more word for telling, explaining and describing at B1.
\end{grammarlens}

\subsection*{Your turn}
\begin{itemize}
\raggedright
  \item \emph{Say it:} \emph{plátnyy}
  \item \emph{Say it:} \emph{plátnyy}, once more
  \item \emph{Say it:} a sentence of your own with \emph{plátnyy}
  \item \emph{Recall:} say the Russian for a payment, then the Russian for free, then say \emph{plátnyy} again
\end{itemize}

\begin{tcolorbox}[breakable,colback=teal!4,colframe=teal!35!black,title={Before you move on}]
What does \ru{платный} mean? (\textquotedblleft{}Paid, that one pays for\textquotedblright{}.) Say it once more.
\end{tcolorbox}
\section[\texorpdfstring{\ru{выгодный} (výgodnyy)}{выгодный (výgodnyy)}]{\ru{выгодный} (výgodnyy) — profitable, a good deal}

\label{lesson:RU-C397-vygodnyy}



\begin{quote}
Before the new one: say the Russian for free, then the Russian for paid.
\end{quote}

\subsection*{You'll want to know: \ru{выгодный}}
\textbf{\ru{выгодный}} (\emph{výgodnyy}) — \textquotedblleft{}profitable, a good deal\textquotedblright{}.

\begin{quote}
\ru{Это} \ru{очень} \ru{выгодное} \ru{предложение}.
\end{quote}

\begin{quote}
\emph{Éto óchen výgodnoye predlozhéniye.}
\end{quote}

\begin{quote}
This is a very good deal.
\end{quote}

\emph{\ru{Выгодно}} — it pays, it is worth it.

\begin{cousinweb}
\textbf{\ru{выгодный}} is built on \emph{\ru{выгода}} \textquotedblleft{}benefit, gain\textquotedblright{}, from \emph{\ru{вы}-} and the root of \emph{\ru{годиться}} \textquotedblleft{}to be fit\textquotedblright{} and \emph{\ru{год}}, Proto-Slavic \emph{*godŭ} \textquotedblleft{}a fitting time\textquotedblright{}: what suits you well.
\end{cousinweb}

\begin{grammarlens}[title={What you've built}]
One more word for telling, explaining and describing at B1.
\end{grammarlens}

\subsection*{Your turn}
\begin{itemize}
\raggedright
  \item \emph{Say it:} \emph{výgodnyy}
  \item \emph{Say it:} \emph{výgodnyy}, once more
  \item \emph{Say it:} a sentence of your own with \emph{výgodnyy}
  \item \emph{Recall:} say the Russian for free, then the Russian for paid, then say \emph{výgodnyy} again
\end{itemize}

\begin{tcolorbox}[breakable,colback=teal!4,colframe=teal!35!black,title={Before you move on}]
What does \ru{выгодный} mean? (\textquotedblleft{}Profitable, a good deal\textquotedblright{}.) Say it once more.
\end{tcolorbox}
